\documentclass[11pt]{article}

\usepackage{amsmath,amssymb}

\usepackage[margin=24mm]{geometry}

\usepackage{longtable}

\allowdisplaybreaks

\newcommand{\Sfun}{\mathcal{S}}

\newcommand{\Cfun}{\mathcal{C}}

\newcommand{\Tfun}{\mathcal{T}}

\title{SoftArm COSSERAT\_PCS Model (1 Section)}

\author{Generated by SoftArm Toolbox}

\date{}

\begin{document}

\maketitle

\section{Model Summary}
This document is generated from the same SymPy model used for numerical code generation.
\begin{description}
\item[Source configuration] \texttt{cosserat\_pcs\_distributed\_tendon\_n1.toml}
\item[Arm family] \texttt{cosserat\_pcs}
\item[Number of sections] 1
\item[Base mode] \texttt{fixed}
\item[Arm inertia model] \texttt{distributed}
\item[Material integration] \texttt{gauss (order 2)}
\end{description}

\section{Notation and Parameters}
The arm base is fixed, hence $\boldsymbol q_B$ is empty.
\begin{equation}
\boldsymbol q_a=\begin{bmatrix}\delta\kappa_{x,1}\\\delta\kappa_{y,1}\\\delta\kappa_{z,1}\\\delta\nu_{x,1}\\\delta\nu_{y,1}\\\delta\nu_{z,1}\end{bmatrix}
\end{equation}

\begin{equation}
\dot{\boldsymbol q}_a=\begin{bmatrix}\dot{\delta\kappa_{x,1}}\\\dot{\delta\kappa_{y,1}}\\\dot{\delta\kappa_{z,1}}\\\dot{\delta\nu_{x,1}}\\\dot{\delta\nu_{y,1}}\\\dot{\delta\nu_{z,1}}\end{bmatrix}
\end{equation}

\begin{equation}
\boldsymbol q=\begin{bmatrix}\boldsymbol q_B\\\boldsymbol q_a\end{bmatrix},\qquad\dot{\boldsymbol q}=\begin{bmatrix}\dot{\boldsymbol q}_B\\\dot{\boldsymbol q}_a\end{bmatrix}
\end{equation}

The NED world frame is used; gravity acts in the positive world $z$ direction.
\subsection{Runtime Parameters}
\begin{longtable}{lll}
\hline
Symbol & Code name & Default \\
\hline
\endfirsthead
\hline Symbol & Code name & Default \\
\hline
\endhead
$L_{1}$ & \texttt{s1\_length} & $0.45$ \\
$m_{1}$ & \texttt{s1\_mass} & $0.18$ \\
$I_{xx,1}$ & \texttt{s1\_Ixx} & $0.0018$ \\
$I_{yy,1}$ & \texttt{s1\_Iyy} & $0.0018$ \\
$I_{zz,1}$ & \texttt{s1\_Izz} & $0.0001$ \\
$\kappa_0_{x,1}$ & \texttt{s1\_kappa0\_x} & $0$ \\
$\kappa_0_{y,1}$ & \texttt{s1\_kappa0\_y} & $0$ \\
$\kappa_0_{z,1}$ & \texttt{s1\_kappa0\_z} & $0$ \\
$\nu_0_{x,1}$ & \texttt{s1\_nu0\_x} & $0$ \\
$\nu_0_{y,1}$ & \texttt{s1\_nu0\_y} & $0$ \\
$\nu_0_{z,1}$ & \texttt{s1\_nu0\_z} & $1$ \\
$EI_{x,1}$ & \texttt{s1\_EI\_x} & $1.1$ \\
$EI_{y,1}$ & \texttt{s1\_EI\_y} & $1.1$ \\
$GJ_{1}$ & \texttt{s1\_GJ} & $0.2$ \\
$GA_{x,1}$ & \texttt{s1\_GA\_x} & $20$ \\
$GA_{y,1}$ & \texttt{s1\_GA\_y} & $20$ \\
$EA_{1}$ & \texttt{s1\_EA} & $50$ \\
$d_{\kappa_x,1}$ & \texttt{s1\_d\_kx} & $0.05$ \\
$d_{\kappa_y,1}$ & \texttt{s1\_d\_ky} & $0.05$ \\
$d_{\kappa_z,1}$ & \texttt{s1\_d\_kz} & $0.02$ \\
$d_{\nu_x,1}$ & \texttt{s1\_d\_vx} & $0.1$ \\
$d_{\nu_y,1}$ & \texttt{s1\_d\_vy} & $0.1$ \\
$d_{\nu_z,1}$ & \texttt{s1\_d\_vz} & $0.1$ \\
$g$ & \texttt{gravity} & $9.81$ \\
$m_e$ & \texttt{tip\_mass} & $0.1$ \\
$I_{e,xx}$ & \texttt{tip\_Ixx} & $0.01$ \\
$I_{e,yy}$ & \texttt{tip\_Iyy} & $0.01$ \\
$I_{e,zz}$ & \texttt{tip\_Izz} & $0.01$ \\
$r_{\mathrm{t1},1}$ & \texttt{act\_t1\_s1\_radius} & $0.0275$ \\
$r_{\mathrm{t2},1}$ & \texttt{act\_t2\_s1\_radius} & $0.0275$ \\
$r_{\mathrm{t3},1}$ & \texttt{act\_t3\_s1\_radius} & $0.0275$ \\
\hline
\end{longtable}

\section{Kinematics}
\begin{equation}
H_{WB}=I_4
\end{equation}

\begin{equation}
H_{BA}=\begin{bmatrix}1 & 0 & 0 & 0.0\\0 & 1 & 0 & 0.0\\0 & 0 & 1 & 0.0\\0 & 0 & 0 & 1\end{bmatrix}\label{eq:mount-transform}
\end{equation}

\subsection{Cosserat Piecewise-Constant-Strain Section}
\begin{equation}
\Tfun(z)=\begin{cases}\dfrac{1-\Sfun(z)}{z},&z\ne0\\\dfrac16,&z=0\end{cases}
\end{equation}

The generalized coordinates are increments from the configured stress-free angular and linear strains.
\begin{equation}
\kappa_i=\kappa_{0,i}+\delta\kappa_i,\qquad\nu_i=\nu_{0,i}+\delta\nu_i
\end{equation}

\begin{equation}
\Omega_i(\xi)=L_i\xi\widehat{\kappa_i},\qquadz_i(\xi)=(L_i\xi)^2\kappa_i^T\kappa_i
\end{equation}

\begin{equation}
R_i=I_3+\Sfun(z_i)\Omega_i+\Cfun(z_i)\Omega_i^2,\qquadr_i=\left[I_3+\Cfun(z_i)\Omega_i+\Tfun(z_i)\Omega_i^2\right]L_i\xi\nu_i
\end{equation}

\begin{equation}
H_i(\xi)=\begin{bmatrix}R_i(\xi)&r_i(\xi)\\0&1\end{bmatrix}\label{eq:local-transform}
\end{equation}

\begin{equation}
H_{Wi}(\xi)=H_{WB}H_{BA}\left(\prod_{k=1}^{i-1}H_k(1)\right)H_i(\xi)
\end{equation}

\begin{equation}
J_{v,i}=\frac{\partial r_{Wi}}{\partial\boldsymbol q},\qquadJ_{\omega,i}^{(:,j)}=\operatorname{vex}\!\left(\operatorname{skew}\!\left(\frac{\partial R_{Wi}}{\partial q_j}R_{Wi}^{T}\right)\right)
\end{equation}

\begin{equation}
J_e=\begin{bmatrix}J_{v,e}\\J_{\omega,e}\end{bmatrix}
\end{equation}


\section{Energy and Dynamics}
Each section uses distributed inertia over $\xi\in[0,1]$. The configured 2-point Gauss--Legendre rule is used.
\begin{equation}
M_i=\int_0^1\!\left(m_iJ_{v,i}^{T}J_{v,i}+J_{\omega,i}^{T}R_{Wi}I_iR_{Wi}^{T}J_{\omega,i}\right)d\xi
\end{equation}

\begin{equation}
M_e=m_eJ_{v,e}^{T}J_{v,e}+J_{\omega,e}^{T}R_eI_eR_e^TJ_{\omega,e}
\end{equation}

\begin{equation}
M=\sum_{i=1}^{N}M_i+M_e\label{eq:mass-assembly}
\end{equation}

\begin{equation}
T=\frac12\dot{\boldsymbol q}^{T}M(\boldsymbol q)\dot{\boldsymbol q}
\end{equation}

\begin{equation}
V=V_g+V_{\mathrm{elastic}},\qquad V_g=-g\!\left(\sum_{i=1}^{N}m_i\int_0^1z_i(\xi)\,d\xi+m_ez_e\right)
\end{equation}

\begin{equation}
V_{\mathrm{elastic}}=\frac12\sum_{i=1}^{N}L_i\left(EI_{x,i}\delta\kappa_{x,i}^2+EI_{y,i}\delta\kappa_{y,i}^2+GJ_i\delta\kappa_{z,i}^2+GA_{x,i}\delta\nu_{x,i}^2+GA_{y,i}\delta\nu_{y,i}^2+EA_i\delta\nu_{z,i}^2\right)
\end{equation}

\begin{equation}
D=\operatorname{diag}\!\left(\begin{bmatrix}d_{\kappa_x,1} L_{1}\\d_{\kappa_y,1} L_{1}\\d_{\kappa_z,1} L_{1}\\d_{\nu_x,1} L_{1}\\d_{\nu_y,1} L_{1}\\d_{\nu_z,1} L_{1}\end{bmatrix}\right)
\end{equation}

\begin{equation}
c(\boldsymbol q,\dot{\boldsymbol q})=\frac{\partial(M\dot{\boldsymbol q})}{\partial\boldsymbol q}\dot{\boldsymbol q}-\frac12\left(\frac{\partial(\dot{\boldsymbol q}^{T}M\dot{\boldsymbol q})}{\partial\boldsymbol q}\right)^T
\end{equation}

\begin{equation}
h=c+\nabla_{\boldsymbol q}V+D\dot{\boldsymbol q}
\end{equation}

The vehicle wrench $w_B$ is expressed in the vehicle body frame, whereas the end wrench $w_e$ is expressed in the NED world frame.
\begin{equation}
Q=S_a\tau_a+B_v(\boldsymbol q)w_B+J_e^Tw_e,\qquadB_v=J_B^T\operatorname{diag}(R_{WB},R_{WB})
\end{equation}

\begin{equation}
M\ddot{\boldsymbol q}+h=Q\label{eq:dynamics}
\end{equation}

\begin{equation}
\dot x=\begin{bmatrix}\dot{\boldsymbol q}\\M^{-1}(Q-h)\end{bmatrix},\qquadx=\begin{bmatrix}\boldsymbol q\\\dot{\boldsymbol q}\end{bmatrix}
\end{equation}


\section{Actuation}
The configured actuator family is \texttt{tendon} with 3 ordered channels.
\begin{center}
\begin{tabular}{ll}
\hline
Channel & Type \\
\hline
\texttt{t1} & \texttt{unilateral} \\
\texttt{t2} & \texttt{unilateral} \\
\texttt{t3} & \texttt{unilateral} \\
\hline
\end{tabular}
\end{center}
\begin{equation}
y(\boldsymbol q_a)=\begin{bmatrix}L_{1} \sqrt{\left(\nu_0_{x,1} + \delta\nu_{x,1}\right)^{2} + \left(- r_{\mathrm{t1},1} \left(\delta\kappa_{y,1} + \kappa_0_{y,1}\right) + \nu_0_{z,1} + \delta\nu_{z,1}\right)^{2} + \left(r_{\mathrm{t1},1} \left(\delta\kappa_{z,1} + \kappa_0_{z,1}\right) + \nu_0_{y,1} + \delta\nu_{y,1}\right)^{2}}\\L_{1} \sqrt{\left(- 0.86602540378443874 r_{\mathrm{t2},1} \left(\delta\kappa_{z,1} + \kappa_0_{z,1}\right) + \nu_0_{x,1} + \delta\nu_{x,1}\right)^{2} + \left(- 0.49999999999999983 r_{\mathrm{t2},1} \left(\delta\kappa_{z,1} + \kappa_0_{z,1}\right) + \nu_0_{y,1} + \delta\nu_{y,1}\right)^{2} + \left(0.86602540378443874 r_{\mathrm{t2},1} \left(\delta\kappa_{x,1} + \kappa_0_{x,1}\right) + 0.49999999999999983 r_{\mathrm{t2},1} \left(\delta\kappa_{y,1} + \kappa_0_{y,1}\right) + \nu_0_{z,1} + \delta\nu_{z,1}\right)^{2}}\\L_{1} \sqrt{\left(- 0.50000000000000042 r_{\mathrm{t3},1} \left(\delta\kappa_{z,1} + \kappa_0_{z,1}\right) + \nu_0_{y,1} + \delta\nu_{y,1}\right)^{2} + \left(0.8660254037844384 r_{\mathrm{t3},1} \left(\delta\kappa_{z,1} + \kappa_0_{z,1}\right) + \nu_0_{x,1} + \delta\nu_{x,1}\right)^{2} + \left(- 0.8660254037844384 r_{\mathrm{t3},1} \left(\delta\kappa_{x,1} + \kappa_0_{x,1}\right) + 0.50000000000000042 r_{\mathrm{t3},1} \left(\delta\kappa_{y,1} + \kappa_0_{y,1}\right) + \nu_0_{z,1} + \delta\nu_{z,1}\right)^{2}}\end{bmatrix}
\end{equation}

\begin{equation}
J_a=\frac{\partial y}{\partial\boldsymbol q_a},\qquad\gamma_a=\dot J_a\dot{\boldsymbol q}_a
\end{equation}

The actuator coordinates are tendon lengths, and positive $T$ denotes tendon tension.
\begin{equation}
\tau_a=-J_a^TT
\end{equation}

Strict actuator-coordinate acceleration is enforced with the full-system Jacobian $J_{a,f}=[0\;J_a]$.
\begin{equation}
\begin{bmatrix}M&J_{a,f}^{T}\\J_{a,f}&0\end{bmatrix}\begin{bmatrix}\ddot{\boldsymbol q}\\T\end{bmatrix}=\begin{bmatrix}Q-h\\\ddot y_{\mathrm{cmd}}-\gamma_a\end{bmatrix}
\end{equation}


\end{document}
